\section*{FTL Symbol Glossary}

This table normalizes only the vocabulary needed to read the present theory manuscript.

\subsection*{Status Labels}

\begin{tabular}{L{0.16\linewidth}L{0.78\linewidth}}
\textbf{Label} & \textbf{Meaning} \\
`[I]` & Imported framework statement fixed before the route analysis. \\
`[P]` & Theorem-level statement proved in the manuscript from its stated hypotheses. \\
`[E]` & Engineering or pre-engineering consequence derived after theorem closure. \\
Soft FTL & Observer-level remote readout regime without ordinary signal propagation in observer space. \\
Hard FTL & Transport regime in which a carrier sector outruns the chosen luminal reference sector. \\
I0 & First integrated architecture in which a soft scan configures a hard entry schedule. \\
\end{tabular}

\subsection*{Core Theory Vocabulary}

\begin{tabular}{L{0.22\linewidth}L{0.72\linewidth}}
\textbf{Term or symbol} & \textbf{Normalized meaning} \\
$L$ & Reference luminal transport sector defining the effective light-cone speed in a chosen observer window. \\
$T$ & Transport or carrier sector compared against $L$ in sector-relative FTL claims. \\
$c_L$ & Effective cone speed of the reference luminal sector. \\
$v_O(T)$ & Observer-relative speed of transport sector $T$. \\
$\mathrm{FTL}_O(T;L)$ & Predicate asserting that sector $T$ is faster than the chosen luminal reference sector for observer $O$. \\
$\Omega_{AB}$ & Candidate soft link-sector shared between two observer windows. \\
Branch-locked sector & Preferred soft-FTL candidate in which local and remote windows read one shared non-factorizable branch structure. \\
Shared-memory bridge & Weak soft precursor based on persistent distributed memory rather than a genuine shared branch structure. \\
Carrier cycle & Hard-FTL entry, hold, and exit cycle on a superluminal transport branch. \\
$B_T$ & Effective carrier-branch quantity used in the hard-transport discussion and benchmarks. \\
$R_{\mathrm{soft}}$ & Soft readout error metric used in the engineering prototype criteria. \\
$\Delta_{\min}$ & Minimum readable margin separating valid soft output classes. \\
\end{tabular}

\subsection*{Canonical FTL Statement Ledger}

\begin{tabular}{L{0.18\linewidth}L{0.14\linewidth}L{0.6\linewidth}}
\textbf{Item} & \textbf{Status} & \textbf{Canonical interpretation} \\
Imported observer-relative semantics & `[I]` & Causal primacy, cone relativity, observer dependence, and admissibility tests govern all soft and hard FTL claims. \\
Branch-locked witness package & `[P]` & First soft theorem package built around binary distinguishability in a shared branch structure. \\
Explicit hard carrier chain & `[P]` & Explicit one-pole and composite hard transport theorems proved in the manuscript. \\
General hard closure chain H-C1--H-C5 & `[P]` & Carrier existence, projected-physical observability, identity preservation, composite binding, observability, and identity compatibility, together with robustness lifting, close the hard route as a theorem package under the explicit perturbation, openness, continuity, and projection-compatibility hypotheses fixed in the hard chapter. \\
Soft/hard split & Core architecture statement & Soft FTL concerns readout; hard FTL concerns transport; they must not be collapsed into one claim. \\
Integrated architecture I0 & `[E]` design target & Soft scan configures hard initialization in a first coupled system. \\
\end{tabular}